\section*{Capítulo \ref{ch:b1:gene-expression} --- Expresión génica: transcripción y traducción}

\begin{solution}{exo:b1:gene-expression:1}
\omterm{def:b1:nucleic-acids:chain}{ARN} $5'$-AUGCCUAAG-$3'$; hebra codificante $5'$-ATGCCTAAG-$3'$.
\end{solution}

\begin{solution}{exo:b1:gene-expression:2}
Una caperuza $5'$ (guanina modificada), una cola $3'$ de poli-A y la
retirada de los intrones por \omterm{prop:b1:gene-expression:processing}{corte y empalme}.
\end{solution}

\begin{solution}{exo:b1:gene-expression:3}
AUG CCG AAA GUU UGA: Met-Pro-Lys-Val, y después parada.
\end{solution}

\begin{solution}{exo:b1:gene-expression:4}
A: se admite y se comprueba el ARNt cargado entrante. P: el ARNt que lleva
la cadena en crecimiento; el enlace se forma entre esa cadena y el
\omterm{def:b1:proteins:aminoacid}{aminoácido} del sitio A. E: el ARNt vaciado, de salida.
\end{solution}

\begin{solution}{exo:b1:gene-expression:5}
Codificantes: $2400 - 600 = 1800$ \omterm{def:b1:nucleic-acids:nucleotide}{nucleótidos}: 599 residuos (600 codones,
uno de parada). Tiempo: $599/4 = \qty{150}{s}$. \omterm{def:b1:gene-expression:ribosome}{Ribosomas}: $1800/90 = 20$ a
la vez.
\end{solution}

\begin{solution}{exo:b1:gene-expression:6}
El mensaje se lee de tres en tres desde un inicio fijo; una o dos bases de
más desplazan todos los codones siguientes y el resto de la \omterm{def:b1:proteins:peptide}{proteína} es
disparate, que suele terminar en una parada prematura. Tres bases de más
añaden un \omterm{def:b1:gene-expression:code}{codón} y restablecen el marco: la \omterm{def:b1:proteins:peptide}{proteína} tiene un residuo de más
y unos pocos equivocados entre las inserciones, y a menudo sigue
funcionando.
\end{solution}

\begin{solution}{exo:b1:gene-expression:7}
CAG (Gln) a UAG (parada): la cadena termina en el residuo 150, una \omterm{def:b1:proteins:peptide}{proteína}
truncada y casi con seguridad inactiva (mutación sin sentido). CAG a CAA:
sigue siendo Gln, silenciosa. CAG a CGG: Arg en lugar de Gln, una sustitución
(de sentido erróneo) cuyo efecto depende de la posición.
\end{solution}

\begin{solution}{exo:b1:gene-expression:8}
El \omterm{def:b1:gene-expression:ribosome}{ribosoma} solo comprueba el apareamiento entre \omterm{def:b1:gene-expression:code}{codón} y \omterm{def:b1:gene-expression:trna}{anticodón}; no puede
ver el \omterm{def:b1:proteins:aminoacid}{aminoácido}. Si una cisteína unida a su ARNt se convierte
químicamente en alanina, el \omterm{def:b1:gene-expression:ribosome}{ribosoma} inserta alanina allí donde el mensaje
dice cisteína: se lee el ARNt, no el \omterm{def:b1:proteins:aminoacid}{aminoácido}. Por tanto, la verdadera
traductora es la sintetasa, que empareja cada \omterm{def:b1:proteins:aminoacid}{aminoácido} con el ARNt
adecuado.
\end{solution}

\begin{solution}{exo:b1:gene-expression:9}
$400\times 4 = 1600$ enlaces; los dos GTP por residuo del \omterm{def:b1:gene-expression:ribosome}{ribosoma} son la
mitad, y los dos de las sintetasas la otra mitad.
\end{solution}

\begin{solution}{exo:b1:gene-expression:10}
En los eucariotas el transcrito se fabrica en el núcleo y los \omterm{def:b1:gene-expression:ribosome}{ribosomas}
están en el \omterm{def:b1:eukaryotic-cell:organelle}{citosol}, separados por la envoltura; y el transcrito no es un
mensajero hasta que se ha empalmado y se le ha puesto la caperuza. La
separación hace posible la maduración misma: el empalme alternativo, el
control de la exportación y la comprobación de que el mensaje está completo
antes de leerlo.
\end{solution}

\begin{solution}{exo:b1:gene-expression:11}
Cuatro mensajeros: con los dos exones (190 \omterm{def:b1:nucleic-acids:nucleotide}{nucleótidos} añadidos), solo con
el 3 (90), solo con el 4 (100) y con ninguno (0). El marco se conserva si la
longitud añadida es múltiplo de 3: 0 y 90 lo conservan; 100 y 190 lo
desplazan, de modo que los mensajeros con el \omterm{def:b1:genomes:gene}{exón} 4 solo o con ambos exones
leen el \omterm{def:b1:genomes:gene}{exón} 5 fuera de marco y truncan la \omterm{def:b1:proteins:peptide}{proteína}. El empalme alternativo
debe respetar el marco, y por esa razón los exones son a menudo múltiplos de
tres.
\end{solution}

\begin{solution}{exo:b1:gene-expression:12}
Arbitrario: nada en la química dice que UUU deba significar Phe; las
asignaciones son convenios que sostienen las sintetasas. Optimizado: la
tercera posición es la más degenerada, de modo que los errores y las
mutaciones que caen ahí son a menudo silenciosos, y los codones que difieren
en una base tienden a codificar \omterm{def:b1:proteins:aminoacid}{aminoácidos} parecidos, con lo que una
sustitución suele ser leve: el código minimiza el daño del error. Congelado:
una sintetasa que cambiara su especificidad alteraría todas las \omterm{def:b1:proteins:peptide}{proteínas} de
la \omterm{def:b1:cell-unit-of-life:cell}{célula} a la vez, miles de ellas, en el mismo instante; casi ninguna
\omterm{def:b1:cell-unit-of-life:cell}{célula} sobrevive a eso, así que el código no puede derivar y está fijado
desde el antepasado común de todos los seres vivos.
\end{solution}

\begin{solution}{pb:b1:gene-expression:1}
\textbf{1.} $30\,000/30 = \qty{1000}{s}$, unos \qty{17}{min}.
\textbf{2.} $28\,000/30\,000 = \qty{93}{\%}$.
\textbf{3.} $30\,000/2000 = 15$.
\textbf{4.} $2\times 30\,000 = \num{60000}$ \omterm{def:b1:water-small-molecules:atp}{ATP}.
\textbf{5.} $1000/10 = 100$ polimerasas sobre el \omterm{def:b1:genomes:gene}{gen}; 360 mensajeros por
hora.
\textbf{6.} Siete intrones, de longitud media $28\,000/7 = \qty{4000}{}$
\omterm{def:b1:nucleic-acids:nucleotide}{nucleótidos}.
\textbf{7.} La transcripción (\qty{17}{min}); los intrones se empalman a
medida que se fabrican, y el último añade solo un minuto.
\textbf{8.} \omterm{def:b1:genomes:gene}{Gen}, $30\,000\times 0.34\,\mathrm{nm} = \qty{10}{\micro m}$; mensajero, \qty{0.68}{\micro m}.
\textbf{9.} $500/4 = \qty{125}{s}$.
\textbf{10.} $2000/80 = 25$.
\textbf{11.} Una cadena termina cada $80/4 = \qty{20}{s}$: 180 por hora.
\textbf{12.} $2/\ln 2 = \qty{2.9}{h}$ de producción plena: unas 520
\omterm{def:b1:proteins:peptide}{proteínas}.
\textbf{13.} $4\times 500 = 2000$ \omterm{def:b1:water-small-molecules:atp}{ATP}.
\textbf{14.} $60\,000/520 = 115$ \omterm{def:b1:water-small-molecules:atp}{ATP} por \omterm{def:b1:proteins:peptide}{proteína}.
\textbf{15.} Unos 2100 \omterm{def:b1:water-small-molecules:atp}{ATP}, el \qty{95}{\%} de ellos en la traducción.
\textbf{16.} $1 - (1 - 10^{-5})^{1503} \approx 1503\times 10^{-5} =
\qty{1.5}{\%}$.
\textbf{17.} $1 - (1 - 10^{-4})^{500} \approx \qty{5}{\%}$.
\textbf{18.} Transcripción: $1.5\%\times 0.75\times 0.67 \approx
0.75\%$ de mensajeros defectuosos y, por tanto, de \omterm{def:b1:proteins:peptide}{proteínas}; traducción:
$5\%\times 0.67 = 3.3\%$; alrededor del \qty{4}{\%} de las moléculas.
\textbf{19.} Un error de \omterm{def:b1:proteins:peptide}{proteína} queda confinado a una molécula, que pronto
se degrada; un error de mensajero se copia en cientos de \omterm{def:b1:proteins:peptide}{proteínas}; un error
de replicación lo heredan todos los descendientes para siempre. El coste de
un error, y por tanto la precisión por la que vale la pena pagar, sube en
cada paso hacia atrás.
\textbf{20.} $3000/180 = 17$ mensajeros presentes.
\textbf{21.} $17\times 25 = 420$ \omterm{def:b1:gene-expression:ribosome}{ribosomas}.
\textbf{22.} $\num{3e9}\times 4 = \num{1.2e10}$ \omterm{def:b1:water-small-molecules:atp}{ATP} al día
$= \qty{2e-14}{mol}$, \qty{1e-9}{J} al día, es decir, \qty{1.2e-14}{W}:
\qty{6}{fW} por micrómetro cúbico.
\textbf{23.} $10^9$ \omterm{def:b1:water-small-molecules:atp}{ATP} por segundo son $\num{8.6e13}$ al día: la síntesis
de \omterm{def:b1:proteins:peptide}{proteínas} es la centésima parte en esta \omterm{def:b1:cell-unit-of-life:cell}{célula} de renovación lenta (en
una bacteria en división es la mitad).
\textbf{24.} No hay mensajero maduro: los transcritos primarios se acumulan
en el núcleo y la \omterm{def:b1:proteins:peptide}{proteína} desaparece a medida que sus mensajeros se
degradan, en cuestión de horas. La \omterm{def:b1:proteins:peptide}{proteína} bacteriana, cuyo \omterm{def:b1:genomes:gene}{gen} no tiene
intrones, no se ve afectada.
\textbf{25.} Unos 2100 \omterm{def:b1:water-small-molecules:atp}{ATP} por molécula: 2000 para la traducción y un
centenar por su parte del transcrito; la primera \omterm{def:b1:proteins:peptide}{proteína}, tras
\qty{1000}{s} de transcripción, un minuto de maduración y exportación y
\qty{125}{s} de traducción: unos veinte minutos.
\end{solution}
